%=====================================================================
%  Chapter 5 — Software Architecture and ROS 2 Integration
%=====================================================================
\chapter{Software Architecture and ROS\,2 Integration}
\label{chap:software}

\section{Introduction}
The previous chapters described the physical robot and the firmware that
animates it. This chapter presents the software that orchestrates the
whole system on the host computer. Because this project is carried out
within a software-engineering curriculum, this is the central chapter of
the report: it explains what ROS\,2 is and why it was adopted, how the
code is organised into packages, how the robot is modelled and
controlled through \texttt{ros2\_control}, how it is simulated in Gazebo,
and how all the pieces connect into a single, coherent system that drives
both the simulation and the real robot.

\section{What ROS\,2 Is, and Why It Is Used Here}
\label{sec:why-ros2}
Despite its name, ROS\,2 (Robot Operating System~2) is not an operating
system but a \emph{middleware} and a set of conventions for building robot
software as a distributed system. Rather than one large program, an
application is structured as a graph of small, independent processes
called \textbf{nodes}, each responsible for a single concern. Nodes never
call each other directly; they communicate through named channels using
three complementary patterns:

\begin{itemize}[itemsep=2pt,topsep=4pt]
  \item \textbf{Topics} --- anonymous publish/subscribe streams, used for
        continuous data such as the inertial measurements or the joint
        states;
  \item \textbf{Services} --- synchronous request/response calls, used for
        short queries such as asking the robot for its status;
  \item \textbf{Actions} --- long-running, cancellable goals that provide
        progress feedback, used for goal-driven motions such as walking or
        bowing.
\end{itemize}

Underneath, ROS\,2 relies on the DDS standard for decentralised discovery
and for configurable quality-of-service, so nodes find one another
automatically and can tune reliability per channel. The source code is
organised into \textbf{packages} gathered in a \textbf{workspace} and
built with the \texttt{colcon} tool.

This model is precisely what makes ROS\,2 the right foundation for Abbes.
First, it imposes a clean \emph{modular} architecture: inverse
kinematics, balance, gait generation, inertial filtering and the
high-level command interface each live in their own node and can be
developed, tested and replaced in isolation. Second, the standardised
message types and the rich tooling (live topic inspection, graph
visualisation, recording and replay, the RViz 3-D viewer) make the system
observable, which is invaluable when debugging a robot. Third, and most
importantly for this project, the \emph{same} interfaces drive both the
Gazebo simulation and the physical robot: an algorithm validated against
the digital twin can be deployed on hardware without rewriting it. Fourth,
\texttt{ros2\_control} provides a standard hardware-abstraction layer
(Section~\ref{sec:ros2control}), and micro-ROS (Chapter~\ref{chap:firmware})
extends the very same graph down to the microcontroller. A monolithic
control program would offer none of this: it would be hard to test in
parts, opaque at runtime, and impossible to retarget from simulation to
hardware without major rewrites.

\section{Workspace and Package Architecture}
\label{sec:packages}
The software is organised as a ROS\,2 workspace whose responsibilities are
split across five packages, each with a clear role:

\begin{itemize}[itemsep=3pt,topsep=4pt]
  \item \texttt{my\_robot\_description} --- the \emph{model} layer. A purely
        static package describing the robot's physical reality (kinematic
        chain, joint axes and limits, link masses, collision geometry,
        visual meshes) in URDF/Xacro. Every other package depends on it.
  \item \texttt{my\_robot\_bringup} --- the \emph{launch} layer. It contains
        no algorithm; its role is orchestration: load the model, start the
        Gazebo simulation, configure the controllers, and bridge the
        simulator to ROS\,2.
  \item \texttt{humanoid\_robot} --- the \emph{algorithm} layer. The nodes
        that actually compute: inertial filtering, balance control, inverse
        kinematics, and the walking-pattern generator
        (Chapter~\ref{chap:kinematics}).
  \item \texttt{humanoid\_msgs} --- the \emph{interface} layer. The action
        and service definitions (\texttt{Walk}, \texttt{StandStill},
        \texttt{SitDown}, \texttt{WaveHand}, \texttt{Bow}, \texttt{Stop},
        \texttt{GetStatus}) shared across the system, so that interfaces are
        never re-defined ad hoc.
  \item \texttt{humanoid\_command\_api} --- the \emph{high-level} layer that
        exposes the robot's command vocabulary and hosts the
        natural-language interface (Chapter~\ref{chap:ai}).
\end{itemize}

This separation mirrors the standard ROS\,2 practice of decoupling the
description, the bring-up, the algorithms and the interfaces, so that a
change in one (for example, a new gait algorithm) does not ripple through
the others.

\section{The Robot Model}
\label{sec:model}
The \texttt{my\_robot\_description} package is the single source of truth
for the robot's geometry. The model is written in Xacro --- a macro
language that expands to URDF --- and was exported from the Onshape CAD
assembly, so that the simulated kinematics match the manufactured parts.
It declares the kinematic chain of revolute joints (the eighteen actuated
degrees of freedom plus fixed joints for the structural links), each with
its axis, parent and child links, angular limits, and the inertial and
collision properties needed for a faithful physics simulation. A
companion file defines the visual materials, and an RViz configuration
provides a consistent debugging view. Because every other package depends
on this model, a single joint-name change here must propagate everywhere
--- a constraint revisited in Section~\ref{sec:conventions}.

Crucially, the model also embeds a \texttt{ros2\_control} block declaring,
for each joint, a position \emph{command interface} and position and
velocity \emph{state interfaces}. This block is what connects the abstract
robot description to the control framework described next.

\section{Motion Control with \texttt{ros2\_control}}
\label{sec:ros2control}
Joint motion is handled by \texttt{ros2\_control}, the standard ROS\,2
control framework, rather than by ad-hoc code. Its central component is
the \textbf{controller manager}, a single node that hosts all the
controllers and owns a \textbf{resource manager}. The resource manager
holds, for every joint, a \emph{command interface} (a value to be written,
here a target position) and \emph{state interfaces} (values that can be
read, here position and velocity). Controllers do not touch the hardware
directly; they \emph{claim} the interfaces they need from the resource
manager, which guarantees that no two controllers fight over the same
joint (Figure~\ref{fig:ros2control}).

\begin{figure}[ht]
  \centering
  \includegraphics[width=0.92\textwidth]{arch_ros2control.pdf}
  \caption{Internals of the \texttt{ros2\_control} controller manager: the
           resource manager owns the command and state interfaces; each
           controller claims the interfaces for its joints, and the whole
           loop runs at 100~Hz.}
  \label{fig:ros2control}
\end{figure}

For Abbes, the joints are grouped into three controllers, configured in
\texttt{controllers.yaml} and all of the \emph{forward command} type
(they forward a commanded position to the position interface):
\begin{itemize}[itemsep=2pt,topsep=4pt]
  \item \texttt{leg\_controller} --- the ten leg joints;
  \item \texttt{arm\_controller} --- the six arm joints;
  \item \texttt{head\_controller} --- the two head joints.
\end{itemize}
A fourth controller, the \texttt{joint\_state\_broadcaster}, is read-only:
it reads every joint's state interfaces and publishes them on
\texttt{/joint\_states}. The whole loop runs at 100~Hz: on each cycle the
hardware interface refreshes the state interfaces, the broadcaster
publishes the joint states, each controller reads its command topic and
writes the command interfaces, and the hardware interface applies them.

Controllers are not all active at start-up; they follow a lifecycle
(Figure~\ref{fig:ctrllifecycle}). A \emph{spawner} loads each controller
from its configuration, asks the resource manager to grant the interfaces
it requires, and then activates it, at which point the controller begins
listening on its command topic and participates in the update loop.

\begin{figure}[ht]
  \centering
  \includegraphics[width=0.92\textwidth]{arch_ctrl_lifecycle.pdf}
  \caption{Controller lifecycle (unloaded $\rightarrow$ configured
           $\rightarrow$ active) and the steps a spawner performs to bring
           a controller online.}
  \label{fig:ctrllifecycle}
\end{figure}

\section{Simulation with Gazebo}
\label{sec:gazebo}
Before any motion is attempted on the physical robot, it is validated on a
\emph{digital twin} in Gazebo Harmonic, whose DART physics engine
simulates gravity, contacts and joint dynamics. The link between Gazebo
and \texttt{ros2\_control} is made by the \texttt{gz\_ros2\_control}
plugin: the URDF declares the plugin and a \texttt{GazeboSimSystem}
hardware interface, so that the simulated joints expose exactly the same
command and state interfaces as a real robot would
(Figure~\ref{fig:gzros2control}). As a result, the controllers described
above are completely unaware of whether they are driving the simulator or
hardware --- which is the foundation of the sim-to-real strategy.

\begin{figure}[ht]
  \centering
  \includegraphics[width=0.6\textwidth]{arch_gz_ros2control.pdf}
  \caption{Integration of Gazebo and \texttt{ros2\_control}: the
           \texttt{gz\_ros2\_control} plugin connects the simulated joints
           to the same interfaces used on hardware, while
           \texttt{ros\_gz\_bridge} relays sensor topics into ROS\,2.}
  \label{fig:gzros2control}
\end{figure}

Internally, Gazebo represents the world as an entity-component system, and
the \texttt{gz\_ros2\_control} plugin runs as part of the simulation loop:
at start-up it parses the \texttt{ros2\_control} block of the URDF and
binds each declared joint to its simulated entity; on every simulation
step it reads the joint states into the state interfaces, lets the
controllers run, and writes the resulting commands back to the simulated
joints (Figure~\ref{fig:gzplugin}). The inertial sensor defined on the
torso is exported to ROS\,2 by the \texttt{ros\_gz\_bridge} as the
\texttt{/imu} topic, mirroring the topic the firmware publishes on the
real robot.

\begin{figure}[ht]
  \centering
  \includegraphics[width=0.92\textwidth]{arch_gz_plugin.pdf}
  \caption{How the \texttt{gz\_ros2\_control} plugin operates inside the
           Gazebo process: it binds the declared joints to simulation
           entities and exchanges state and commands on every step.}
  \label{fig:gzplugin}
\end{figure}

The \texttt{my\_robot\_bringup} package ties this together in a single
launch sequence: it publishes the robot description (expanding the Xacro),
starts Gazebo with an empty world, spawns the robot from its description,
starts the \texttt{ros\_gz\_bridge} for the clock and IMU, and finally
spawns the four controllers in order.

\section{The Node and Topic Graph}
\label{sec:graph}
Putting the pieces together yields the system graph of
Figure~\ref{fig:nodegraph}. At the top, the algorithm nodes of the
planning layer (inverse kinematics, balance, and the gait generator of
Chapter~\ref{chap:kinematics}) publish target joint positions, in radians,
onto the three controller command topics. The controller manager applies
them through the controllers described above. From there, the graph splits
into two paths that share the same upstream:

\begin{itemize}[itemsep=2pt,topsep=4pt]
  \item the \textbf{simulation path}, where the \texttt{gz\_ros2\_control}
        plugin drives the joints inside Gazebo's DART physics and the
        simulated IMU is bridged back to \texttt{/imu};
  \item the \textbf{real-robot path}, where the commands are forwarded to
        the ESP32 over micro-ROS and the MPU6050 publishes \texttt{/imu}.
\end{itemize}

Both paths converge on the same downstream topics --- \texttt{/imu},
\texttt{/joint\_states} and the transform tree \texttt{/tf} --- which the
algorithm nodes consume, closing the loop. Because the real robot listens
on a single command topic in servo degrees while the controllers publish
three topics in radians, a dedicated bridge node converts and re-groups
the commands for the firmware; this sim-to-real bridge is what allows the
poses validated in simulation to be replayed on the physical robot.

\begin{figure}[ht]
  \centering
  \includegraphics[width=\textwidth]{arch_node_graph.pdf}
  \caption{The complete node and topic graph: the planning layer feeds the
           controller manager, which drives either the Gazebo simulation
           path or the real-robot (ESP32) path; both share the downstream
           sensor and state topics.}
  \label{fig:nodegraph}
\end{figure}

\section{Cross-Cutting Conventions}
\label{sec:conventions}
A consistent system depends on a few conventions enforced across all
packages:
\begin{itemize}[itemsep=2pt,topsep=4pt]
  \item \textbf{Units.} The entire ROS\,2 side works in \emph{radians},
        following the URDF zero convention; the conversion to servo degrees
        happens only at the firmware/bridge boundary
        (Chapter~\ref{chap:firmware}), never inside the ROS\,2 nodes.
  \item \textbf{Joint ordering.} The order of the joints must match
        \texttt{controllers.yaml} exactly; a single mismatch would silently
        send a command to the wrong joint.
  \item \textbf{Interfaces.} All message, service and action types live in
        \texttt{humanoid\_msgs}, never defined ad hoc in an implementation
        package.
  \item \textbf{Actions vs. topics.} Goal-driven motions use action servers
        (cancellable, with feedback); streaming data such as the IMU and
        joint states use topics.
\end{itemize}

\section{Conclusion}
This chapter presented the software architecture of Abbes: a modular
ROS\,2 system split into description, bring-up, algorithm, interface and
command packages; a robot model that feeds \texttt{ros2\_control}; a Gazebo
digital twin that exposes the very same interfaces as the hardware; and a
node graph whose simulation and real-robot paths share a common upstream
and a common set of sensor topics. This architecture is what makes the
project tractable as software --- observable, modular, and portable from
simulation to hardware. The next chapter looks inside the algorithm layer,
at the kinematics and gait-planning that compute the joint targets flowing
through this graph.
